% Simple overview deck for the 2D-turbulence PIV data-generation project.
% Self-contained: compiles on Overleaf with pdfLaTeX, no image uploads needed.
% (To use on Overleaf: new Beamer project, replace main.tex with this file.)

\documentclass[aspectratio=169]{beamer}

\usetheme{Madrid}
\usecolortheme{seahorse}
\setbeamertemplate{navigation symbols}{}

\title{Synthetic PIV Training Data from 2D Turbulence}
\subtitle{Status \& plan}
\author{Arup Mazumder}
\institute{Surf Lab, University of Rhode Island}
\date{June 2026}

\begin{document}

\begin{frame}
  \titlepage
\end{frame}

% ---------- Slide 1: goal + pipeline ----------
\begin{frame}{Goal \& pipeline}
\textbf{Goal.} Generate synthetic Particle Image Velocimetry (PIV) image pairs with
\emph{known} velocity fields, to train an ML model on turbulent flow.

\vspace{1em}
\textbf{Pipeline (pure Julia):}
\begin{itemize}\setlength\itemsep{0.4em}
  \item \texttt{2DTurbulence.jl} --- Oceananigans 2D simulation, $512\times512$ periodic
        grid, Lagrangian tracer particles.
  \item \texttt{CombineAndConquer.jl} --- merges field + particle outputs.
  \item \texttt{ImageGen.jl} --- renders particle images at frames $A$, $B$ and saves the
        ground-truth velocity field as the label.
\end{itemize}
\end{frame}

% ---------- Slide 2: what we measured ----------
\begin{frame}{What we measured about the simulation}
\begin{itemize}\setlength\itemsep{0.6em}
  \item \textbf{Displacement floor} $s_{\max}\approx 5$ px --- the smallest motion between
        two saved frames. Sets which displacement bins are achievable.
  \item \textbf{Decorrelation time} $\tau \approx 30$ time units --- the flow is essentially
        \emph{frozen} over the short interval an image pair spans.
\end{itemize}

\vspace{1em}
\begin{block}{Key consequence}
\centering
\textbf{One simulation $\approx$ one independent flow label.}\\
Extra frames are near-duplicates --- diversity must come from \emph{many simulations}.
\end{block}
\end{frame}

% ---------- Slide 3: design decisions ----------
\begin{frame}{Design decisions for scaling up}
\begin{itemize}\setlength\itemsep{0.45em}
  \item \textbf{Many simulations, one image pair each} --- diversity from varied random
        seeds, not from many frames per run.
  \item \textbf{When to sample:} at the eddy-turnover time $t = C/(U\,k_p)$, where $k_p$ is
        the energy-weighted wavenumber --- samples every run at the same dynamical age.
  \item \textbf{Storage:} keep only velocity fields $(u,v)$ + metadata; vorticity, speed,
        etc.\ are recomputed on demand ($\sim$60\% smaller, faster).
  \item \textbf{Metadata per sample:} seed, all parameters, git commit, spectrum, shear,
        vorticity --- everything needed to regenerate the exact dataset.
\end{itemize}
\end{frame}

% ---------- Slide 4: status + next ----------
\begin{frame}{Status \& next steps}
\textbf{Done}
\begin{itemize}\setlength\itemsep{0.2em}
  \item Pipeline verified end-to-end (Mac Studio, CPU); displacement bins $[10,20,30]$ px.
  \item Pipeline bugs fixed (reproducible \texttt{-{}-seed}, output trimmed to $u,v$, dead
        code removed); tracked in \texttt{bugs.md}.
\end{itemize}

\vspace{0.6em}
\textbf{Next}
\begin{itemize}\setlength\itemsep{0.2em}
  \item Generate a $\sim$1000-sample dataset on the Unity cluster (parallel array job).
  \item Train a baseline ML model.
  \item Later: vary initial-condition parameters for \emph{regime} diversity (v2).
\end{itemize}
\end{frame}

\end{document}
